% =====================================================================
%  5 RESULTADOS E DISCUSSÃO — "O que descobri? O que isso significa?"
%  É a essência do trabalho e o principal critério de avaliação.
%  As subseções seguem EXATAMENTE as etapas do Método (Dica 6).
%  Para cada etapa: o que foi obtido (tabular os dados) + análise crítica.
%  Resultados negativos e soluções de contorno também são resultados
%  (veja como o relatório do INCRA trata a limitação do Autodesk Tandem).
% =====================================================================
\section{Resultados e Discussão}\label{sec:resultados}

\escrever{parágrafo de abertura: o que esta seção apresenta e como está organizada.}

\subsection{Estado da arte}
\escrever{síntese dos achados da Etapa 1 que orientaram as decisões de projeto.}

\subsection{Fabricação dos sensores}
\escrever{fotos dos protótipos (Figura), parâmetros finais, falhas de impressão e ajustes.}

\subsection{Instrumentação e firmware}
\escrever{leitura do sensor, ruído medido, estabilidade da linha de base.}

\subsection{Caracterização em bancada}
\escrever{curvas força $\times$ $\Delta B$ (Gráfico), histerese, repetibilidade.}

\subsection{Calibração}

A Tabela~\ref{tab:modelos} compara os modelos de calibração.
\escrever{substituir pelos resultados com o sensor fabricado.}

\begin{table}[H]
  \centering
  \caption{Erro de estimativa de força por modelo de calibração}
  \label{tab:modelos}
  \begin{spacing}{1.0}\small
  \begin{tabular}{lccc}
    \toprule
    \textbf{Modelo} & \textbf{RMSE (N)} & \textbf{R$^2$} & \textbf{Parâmetros} \\
    \midrule
    Regressão linear              & -- & -- & 4 \\
    Regressão polinomial (grau 2) & -- & -- & 10 \\
    MLP 2$\times$16               & -- & -- & 353 \\
    MLP 2$\times$128              & -- & -- & 17.153 \\
    \bottomrule
  \end{tabular}
  \fonte{Elaborado pelo autor (2026).}
  \end{spacing}
\end{table}

\escrever{discussão: qual modelo foi escolhido para o ESP32 e por quê, comparando com \citeonline{pattabiraman2025eflesh} e \citeonline{barthel2026eflesh}.}

\subsection{Integração e feedback vibrotátil}
\escrever{interferência do ERM em função da distância (Gráfico), latência da malha, consumo.}

\subsection{Validação}
\escrever{desempenho nas tarefas com objetos reais.}

\subsection{Discussão}
\escrever{o que os resultados significam diante da literatura; limitações encontradas e por quê.}
